\documentclass[10pt,a4paper]{amsart}
\usepackage[margin=19mm]{geometry}
\usepackage{amsmath,amssymb,mathtools}
\usepackage[scr=rsfs,cal=euler]{mathalfa}
\usepackage{xfrac}
\usepackage[unicode,hidelinks,bookmarks=false]{hyperref}
\newcommand{\ie}{i.e.\ }
\newcommand{\cf}{cf.\ }
\newcommand{\resp}{respectively\ }
\newcommand{\Subsection}{\subsection}
\providecommand{\nobreakdash}{\nobreak\hbox{-}}
\setlength{\emergencystretch}{2em}
\multlinegap0pt
\usepackage{mathtools}
\usepackage{amssymb}
\newtheorem{lemma}{Lemma}
\newtheorem{proposition}{Proposition}
\newcommand{\Hom}{\opname{Hom}}
\newcommand{\add}{\opname{add}\nolimits}
\newcommand{\cd}{{\mathcal D}}
\newcommand{\der}{\cd}
\newcommand{\opname}[1]{\operatorname{\mathsf{#1}}}
\newcommand{\rank}{\opname{rank}\nolimits}
\newcommand{\thick}{\opname{thick}}
\title{On endomorphism algebras of silting complexes over hereditary abelian categories}
\author{Wei Dai, Changjian Fu, Liangang Peng}
\date{Source: arXiv:2602.17197v3 (CC BY 4.0)}
\begin{document}
\maketitle

\noindent\emph{Source and licence.} On endomorphism algebras of silting complexes over hereditary abelian categories. Version arXiv:2602.17197v3 (CC BY 4.0), distributed by arXiv under the Creative Commons Attribution 4.0 International licence, http://creativecommons.org/licenses/by/4.0/, as recorded in the arXiv licence metadata for this exact version and shown on the abstract page https://arxiv.org/abs/2602.17197v3. Full licence notice: \texttt{licenses/arxiv-2602.17197-CC-BY-4.0.txt}.\par\smallskip

\noindent\textit{Editorial scope.} Counted unit: the article's proposition ind-perp-completion (source0.tex lines 581--583). The article's complete original proof of it is delivered with it (source0.tex lines 584--611, one \texttt{proof} environment of 28 lines). No mathematics is written, repaired or re-derived: the statement and the proof are the article's own lines, in the article's own notation, and no retained line is substituted or re-flowed.  Retained uncounted as the source-local context the proof needs: lines 511--513; lines 517--525.  Nothing was omitted from any retained span, so the package carries no omission record.  No retained line contains a citation, so the wrapper carries no bibliography and no bibliography entry.  No retained line contains a figure environment, an image-inclusion command or drawing code, so no image is delivered and the package declares no asset dependency.  Editorial formatting only, in addition: an AMS \texttt{amsart} preamble that restates the article's own declarations and macros used by the retained lines, namely the environments \texttt{lemma}, \texttt{proposition} on separate counters, together with the standard packages those lines need.  The retained environments print in this document as Lemma 1, Proposition 1 in order of appearance, whereas the article prints the same items with the numbering of its own full document.  The complete native source of the article as distributed in the arXiv e-print for this exact version is bundled unchanged as \texttt{sources/194946/source0.tex}.
\begin{lemma}\label{l:no-cycle}
Let $T$ be an object in $\der^b(\mathcal{H})$ such that $\Hom_{\der^b(\mathcal{H})}(T, T[1])=0$. Then the subcategory $\add T$ has no cycle.
\end{lemma}

\begin{lemma}\label{l:no-cycle-plus} Let $T$ be an object in $\der^b(\mathcal{H})$ such that $\Hom_{\der^b(\mathcal{H})}(T, T[1])=0$.
\begin{itemize}
    \item[(a)]  There is a unique decomposition $T =  T_1 \oplus T_2 \oplus \cdots \oplus T_r$ into indecomposable summands $T_i \in \mathcal{H}[n_i]$ with integers $n_1 \leq n_2 \leq \cdots \leq n_r$, such that $\Hom_{\der^b(\mathcal{H})}(T_j, T_i) = 0$ for all $i < j$.
    \item[(b)] Furthermore, if $T$ is a presilting object, then 
    \[
\thick T = \thick T_r * \thick T_{r-1} * \cdots * \thick T_1.
\]
\end{itemize}
\end{lemma}

\begin{proposition}\label{p: ind-perp-completion}
Let $N$ be an almost silting object in $\der^b(\mathcal{H})$. Then there exists an indecomposable presilting object $D$, such that $D \oplus N$ is silting and $N \in (\thick D)^\perp$.
\end{proposition}

\begin{proof}
Let $\rank K_0(\mathcal{H})=n$.
Without loss of generality, assume that $N$ is basic. By Lemma \ref{l:no-cycle-plus} (a), we fix a decomposition  $N =  N_1 \oplus N_2 \oplus \cdots \oplus N_{n-1}$ of $N$, where each $N_i$ is an indecomposable summand such that $N_i \in \mathcal{H}[m_i]$ with $m_1 \leq m_2 \leq \cdots \leq m_{n-1}$ in $\mathbb{Z}$, and $\Hom_{\der^b(\mathcal{H})}(N_j, N_i) = 0$ for all $i < j$. 

Let $T=N\oplus X$ be a basic silting object of $\der^b(\mathcal{T})$. Define 
\[
p_T=\max\{l+1 \mid \Hom_{\der^b(\mathcal{H})}(X,N_i[\mathbb{Z}])=0, 1\leq i\leq l\}.
\]
If $p_T=n$, then $N\in (\thick X)^\perp$ as desired.  If $p_T<n$,  it suffices to  show that there is a silting object $M$ containing $N$ as a direct summand such that $p_M>p_T$.

Now assume that $p:=p_T<n$. 
Let 
\begin{equation}\label{tri:left-mutation}
    X\xrightarrow{f} N'\to Y\to X[1]
\end{equation} 
be the triangle associated with the left mutation of $T$ at $X$. By the definition of $p$, we have $N'\in \add N_p\oplus\cdots\oplus N_{n-1}$, which implies $\Hom_{\der^b(\mathcal{H})}(N',N_i[\mathbb{Z}])=0$ for any $1\leq i<p$. For any $1\leq i<p$, applying $\Hom_{\der^b(\mathcal{H})}(-,N_i[\mathbb{Z}])$ to the triangle \eqref{tri:left-mutation} yields an exact sequence
\[
0=\Hom_{\der^b(\mathcal{H})}(X[1],N_i[\mathbb{Z}])\to \Hom_{\der^b(\mathcal{H})}(Y,N_i[\mathbb{Z}])\to \Hom_{\der^b(\mathcal{H})}(N',N_i[\mathbb{Z}])=0.
\]
Consequently, $\Hom_{\der^b(\mathcal{H})}(Y,N_i[\mathbb{Z}])=0$ and hence $p_{\mu_X^-(T)} \geq p$.




If $N_p\in \add N'$, then $\Hom_{\der^b(\mathcal{H})}(X,N_p)\neq 0$ and $\Hom_{\der^b(\mathcal{H})}(N_p,Y)\neq 0$. Since $Y\oplus N_p$ is a presilting object,  $\Hom_{\der^b(\mathcal{H})}(Y,N_p)=0$ by Lemma \ref{l:no-cycle}. Moreover, $\Hom_{\der^b(\mathcal{H})}(Y,N_p[>0])=0$.  Note that $N_p\in \mathcal{H}[m_p]$ and $\Hom_{\der^b(\mathcal{H})}(N_p,Y)\neq 0$, it follows that $Y\in \mathcal{H}[m_p]$ or $\mathcal{H}[m_p+1]$. We conclude that $\Hom_{\der^b(\mathcal{H})}(Y,N_p[<0])=0$, and hence $\Hom_{\der^b(\mathcal{H})}(Y,N_p[\mathbb{Z}])=0$. Therefore $p_{\mu_X^-(T)}\geq p+1>p_T$ and we may take $M=\mu_X^-(T)$.

Now assume that $N_p\not\in \add N'$. Without loss of generality, we assume that $N'\in \add N_k\oplus\cdots\oplus N_{n-1}$ and $N_k\in \add N'$ for some $p<k\leq n-1$. It follows that $\Hom_{\der^b(\mathcal{H})}(X,N_p)=0$ and $\Hom_{\der^b(\mathcal{H})}(X,N_k)\neq 0$.  By the definition of $p_T$, we know that $\Hom_{\der^b(\mathcal{H})}(X,N_p[-u])\neq 0$ for some $u\geq 1$. Consequently, $X\in \mathcal{H}[m_p-u]$ or $\mathcal{H}[m_p-u-1]$. On the other hand, by $\Hom_{\der^b(\mathcal{H})}(X,N_k)\neq 0$, we have $X\in \mathcal{H}[m_k]$ or $\mathcal{H}[m_k-1]$. Noticing that $m_p\leq \cdots\leq m_k$, we conclude that $m_p=\cdots=m_k$, $u=1$ and $X\in \mathcal{H}[m_p-1]$. Consequently, $Y\in \mathcal{H}[m_p]$ and $\Hom_{\der^b(\mathcal{H})}(Y,N_p)\cong \Hom_{\der^b(\mathcal{H})}(X[1],N_p)\neq 0$. Now consider the left mutation  of $N\oplus Y$ at $Y$. By the above discussion, we conclude that $p_{\mu_Y^-(N\oplus Y)}>p$ and we may take $M=\mu_Y^-(N\oplus Y)$. This completes the proof.
\end{proof}

\end{document}
